\documentclass[a4paper,12pt]{article}
\usepackage[T2A]{fontenc}
\usepackage[utf8]{inputenc}
\usepackage[russian]{babel}
\usepackage{amsmath,amssymb,mathtools}
\usepackage{geometry}
\geometry{left=20mm,right=20mm,top=20mm,bottom=23mm}
\renewcommand{\familydefault}{\sfdefault}
\usepackage{microtype}
\usepackage{array,booktabs,tabularx}
\usepackage{tikz}
\usepackage{xcolor}
\usepackage{fancyhdr}
\usepackage[unicode=true,colorlinks=true,linkcolor=darkaccent]{hyperref}
\definecolor{accent}{HTML}{087C83}
\definecolor{darkaccent}{HTML}{07595D}
\definecolor{ink}{HTML}{27333B}
\definecolor{wrong}{HTML}{A53639}
\color{ink}
\setlength{\parindent}{0pt}
\setlength{\parskip}{5pt}
\linespread{1.13}
\pagestyle{fancy}
\fancyhf{}
\renewcommand{\headrulewidth}{0.25pt}
\fancyhead[L]{\scriptsize \textcolor{darkaccent}{МАТЕМАТИКА / РЕВЬЮ ДОМАШНЕЙ РАБОТЫ}}
\fancyhead[R]{\scriptsize 10.11.2026}
\fancyfoot[C]{\scriptsize Лёвин Артём Александрович эксклюзивно для Николь Саркисьянц}\fancyfoot[R]{\scriptsize\thepage}
\newcommand{\errorlabel}{\textcolor{wrong}{\textbf{Ошибка.}}}
\newcommand{\keyidea}{\textcolor{darkaccent}{\textbf{Ключевая идея.}}}
\newcommand{\result}{\textcolor{darkaccent}{\textbf{Итог.}}}
\newcommand{\subline}[1]{\par\vspace{3pt}{\color{darkaccent}\bfseries #1}\par}
\begin{document}
\thispagestyle{empty}
\begin{center}
\vspace*{28mm}
{\color{darkaccent}\large МАТЕМАТИКА}\par\vspace{15mm}
{\Huge Ревью домашней работы}\par\vspace{11mm}
{\fontsize{17}{23}\selectfont\color{accent}Стереометрия}\par\vspace{15mm}
{\small Подготовка к ЕГЭ по профильной математике}\par
\vfill
{\normalsize Николь Саркисьянц}\par\vspace{8mm}
{\normalsize 10 ноября 2026 г.}\par\vspace{12mm}
{\small Лёвин Артём Александрович}\par
{\small эксклюзивно для Николь Саркисьянц}\par\vspace{12mm}
\begin{tikzpicture}\draw[accent,line width=.9pt] (0,0) .. controls (1.1,.35) and (2.1,-.35) .. (3.2,0) .. controls (4.4,.35) and (5.7,-.35) .. (7.0,0);\end{tikzpicture}
\end{center}
\clearpage

\section*{Сводная таблица}
\textcolor{darkaccent}{Источник:} «09.11.26.pdf», тренировочный блок, задача №~8 (фотография 09.jpg).

В таблицу включено задание, в котором правильный численный ответ получен с нарушением формулы площади поверхности.

\small\renewcommand{\arraystretch}{1.4}
\begin{tabularx}{\linewidth}{@{}p{11mm}p{28mm}X p{29mm}p{32mm}@{}}
\toprule
\textbf{№} & \textbf{Тема} & \textbf{Суть ошибки} & \textbf{Ваш ответ} & \textbf{Верный ответ}\\\midrule
\hyperref[task:8]{8} & Призма & Вместо боковой площади подставлен удвоенный объём. & \shortstack[l]{$V=18$\\$S=36+6\sqrt3$} & \shortstack[l]{$V=18$\\$S=36+6\sqrt3$}\\
\bottomrule
\end{tabularx}
\normalsize

\vspace{8mm}
Равенство численных ответов объясняется особенностями заданных длин. Корректность численного результата здесь не подтверждает корректность использованного правила.

\clearpage
\thispagestyle{empty}
\vspace*{\fill}
\begin{center}{\fontsize{25}{30}\selectfont\color{darkaccent}Разбор ошибок}\par\vspace{5mm}
{\large Подробное объяснение и задача на закрепление}\end{center}
\vspace*{\fill}
\clearpage
\phantomsection\label{task:8}
\section*{Задание № 8}

\subline{Условие}
Все рёбра правильной треугольной призмы равны $2\sqrt3$. Найдите её объём и площадь полной поверхности.

\subline{Ваше решение}
Вы обозначили общую длину ребра через $a=2\sqrt3$, верно определили объём $V=18$, периметр основания $P_0=6\sqrt3$ и суммарную площадь двух оснований $2S_0=6\sqrt3$. При вычислении полной поверхности была записана строка
\[
S_{\text{полн}}=2\cdot18+6\sqrt3=36+6\sqrt3.
\]
Итоговые значения объёма и площади поверхности совпадают с правильными.

\textbf{Ваш ответ:} $V=18$, $S_{\text{полн}}=36+6\sqrt3$.

\subline{Где именно возникает ошибка}
\textbf{Последний правильный шаг:} $V=18$, $P_0=6\sqrt3$, $2S_0=6\sqrt3$.

\errorlabel\ При вычислении полной поверхности в качестве боковой площади использовано выражение $2\cdot18=2V$. В Вашем решении число $18$ ранее обозначало объём; переход к удвоенному объёму для нахождения площади не обоснован ни одной формулой.

\subline{Почему этот шаг неверен}
Объём и площадь измеряют разные величины. Объём измеряется в кубических единицах, а площадь --- в квадратных. Удвоение объёма не превращает его в боковую поверхность. Для \emph{прямой} призмы боковые грани являются прямоугольниками. Их общая площадь находится по формуле
\[
S_{\text{бок}}=P_0h,
\]
где $P_0$ --- периметр основания, $h$ --- высота призмы. В правильной треугольной призме высота равна боковому ребру.

В этой конкретной задаче $2V=36$ \emph{численно совпало} с $S_{\text{бок}}=36$. При других размерах призмы это совпадение обычно исчезает.

\keyidea\ Для объёма используйте площадь основания: $V=S_0h$. Для боковой поверхности используйте периметр основания: $S_{\text{бок}}=P_0h$. Для полной поверхности добавьте площади обоих оснований.

\subline{Исправление с последнего правильного шага}
После получения $P_0=6\sqrt3$ и $h=2\sqrt3$ вычислите именно боковую площадь:
\[
S_{\text{бок}}=P_0h=6\sqrt3\cdot2\sqrt3=36.
\]
Тогда
\[
S_{\text{полн}}=S_{\text{бок}}+2S_0=36+6\sqrt3.
\]

\subline{Полное правильное решение}
\textbf{1. Найдём площадь основания.} Основание --- равносторонний треугольник со стороной $a=2\sqrt3$. Поэтому
\[
S_0=\frac{a^2\sqrt3}{4}=\frac{(2\sqrt3)^2\sqrt3}{4}=3\sqrt3.
\]
\textbf{2. Найдём объём.} Высота правильной треугольной призмы совпадает с длиной её бокового ребра: $h=2\sqrt3$. Получаем
\[
V=S_0h=3\sqrt3\cdot2\sqrt3=18.
\]
\textbf{3. Найдём периметр основания и боковую площадь.}
\[
P_0=3a=6\sqrt3,\qquad S_{\text{бок}}=P_0h=36.
\]
\textbf{4. Найдём полную площадь поверхности.}
\[
S_{\text{полн}}=2S_0+S_{\text{бок}}=6\sqrt3+36.
\]

\subline{Проверка}
У призмы три прямоугольные боковые грани. Площадь каждой:
\[
S_{\text{грани}}=a\cdot h=(2\sqrt3)^2=12.
\]
Суммарная площадь трёх боковых граней равна $3\cdot12=36$, что совпадает с результатом по формуле $P_0h$. Площадь двух треугольных оснований равна $2\cdot3\sqrt3=6\sqrt3$. Размерности согласованы: объём выражен в кубических единицах, площадь --- в квадратных.

\result\ Объём $18$, площадь полной поверхности $36+6\sqrt3$. Конечный ответ верен; для площади поверхности требуется исправить способ получения числа $36$.

\subline{Задача на закрепление}
В правильной треугольной призме все рёбра равны $4$. Найдите объём и площадь полной поверхности.
\end{document}
